\chapter{Conclusiones}
A continuación se responden las preguntas propuestas por la cátedra
a partir de las actividades realizadas en el TP2.
 
\begin{enumerate}[leftmargin=*, label=\textbf{\arabic*.}]
 
    \item \textbf{Diferencias entre las curvas I–V del 1N4007 (silicio) y el 1N60 (germanio), y su influencia en la elección del componente}\\
    Las curvas obtenidas en simulación y laboratorio muestran que el 1N60 comienza a conducir a tensiones directas notablemente menores que el 1N4007: mientras el germanio presenta conducción apreciable a partir de valores cercanos a 0,2V, el silicio requiere aproximadamente 0,5V para iniciar la conducción. Las tensiones de umbral medidas en laboratorio confirmaron esto, con caídas de entre 297--355mV para el 1N60 y 567--712mV para el 1N4007. En aplicaciones de baja tensión, donde la caída en el diodo representa una fracción significativa de la tensión disponible, el diodo de germanio es preferible. Sin embargo, el 1N60 presenta mayor corriente de fuga en inversa y mayor sensibilidad térmica, lo que lo hace menos adecuado para entornos con variaciones de temperatura importantes.

    \item \textbf{Desplazamiento de la curva característica con la temperatura y su efecto en la estabilidad del circuito}\\
    Las simulaciones con barrido de temperatura (20°C, 100°C y 125°C) mostraron que al aumentar la temperatura la curva I-V se desplaza hacia la izquierda en polarización directa: para un mismo voltaje aplicado, la corriente que circula es mayor. Esto se debe a que la barrera de potencial de la unión PN disminuye con la temperatura, facilitando el paso de portadores. En polarización inversa, la corriente de fuga aumenta considerablemente con la temperatura, efecto más pronunciado en el 1N60 por sus propiedades intrínsecas como diodo de germanio. En un circuito diseñado para operar en entornos con variaciones térmicas, este comportamiento implica que el punto de operación se desplaza con la temperatura, lo que puede alterar las corrientes y tensiones de trabajo y debe ser contemplado en el diseño mediante márgenes adecuados o compensación térmica.

    \item \textbf{Importancia de respetar VRRM e IFSMI, y consecuencias de superar la tensión de ruptura en un diodo no Zener}\\
    VRRM es la tensión inversa repetitiva máxima que el diodo puede soportar sin dañarse. Superarla de forma sostenida puede provocar la ruptura irreversible de la juntura. IFSMI es la corriente de pico no repetitiva máxima en directa; excederla, aunque sea brevemente, puede destruir el dispositivo por disipación térmica excesiva. En un diodo que no es tipo Zener, superar la tensión de ruptura no produce una limitación controlada de tensión como en el Zener, sino una avalancha descontrolada: la corriente inversa crece abruptamente sin mecanismo de autorregulación, generando una disipación de potencia que destruye el componente de forma permanente si no se limita externamente.

    \item \textbf{Relación entre la tensión nominal VZ y la forma de onda en el osciloscopio, y comportamiento cuando el Zener queda polarizado en directa}\\
    Las formas de onda obtenidas en el osciloscopio confirmaron el comportamiento teórico: la señal de salida quedó recortada a niveles aproximadamente iguales a las tensiones Zener de los diodos utilizados (5,1V y 6,8V). Cuando el diodo Zener queda polarizado en directa durante el semiciclo opuesto, se comporta como un diodo rectificador convencional, imponiendo una caída de aproximadamente 0,7V. Esto explica por qué los niveles de recorte calculados fueron VZ1 + 0,7V en un semiciclo y $-(V_{Z2}+0,7\,V)$ en el otro, tal como se verificó tanto en la simulación como en el laboratorio.

    \item \textbf{Discrepancias entre valores medidos con multímetro y gráficas de simulación, e influencia de las resistencias internas y tolerancias}\\
    Las mediciones de tensión umbral en laboratorio arrojaron valores ligeramente inferiores a los esperados por los modelos teóricos. En el 1N4007 la tensión directa medida fue de aproximadamente 0,57V para corrientes bajas, frente a los 0,7V del modelo simplificado. Esto se explica por dos factores: por un lado, la resistencia interna del voltímetro introduce una pequeña carga paralela que puede modificar levemente la medición; por otro, los modelos SPICE utilizados en la simulación fueron ajustados con parámetros del fabricante que pueden diferir del componente físico real. Las tolerancias de fabricación de los diodos también contribuyen: la corriente de saturación IS, que define el inicio de la conducción, varía entre unidades del mismo modelo. En conjunto, estas fuentes de error son esperables y no comprometen la validez del análisis, ya que las desviaciones observadas son consistentes con las características declaradas en las hojas de datos. 

    \item \textbf{Función de la resistencia limitadora de corriente en los recortadores y criterio para calcular su valor óptimo}\\
    La resistencia en serie en los circuitos recortadores cumple una función doble: limita la corriente que circula por los diodos Zener cuando están en conducción, protegiéndolos de la destrucción por exceso de potencia disipada; y determina la caída de tensión en la rama serie cuando el circuito no está en recorte, afectando la fidelidad de la señal de salida respecto a la entrada. Para calcular su valor óptimo se debe considerar que la corriente máxima por el Zener no supere $I_{ZM}=\dfrac{P_{Zmax}}{V_Z}$, lo que impone una cota inferior: $R \geq \dfrac{V_{in,max}-V_Z}{I_{ZM}}$ . Al mismo tiempo, una resistencia demasiado alta reduce la corriente en conducción por debajo de IZmin, sacando al Zener de su zona de regulación. En los circuitos implementados se utilizó $R=1\,k\Omega$, valor que resultó adecuado para la señal de 20V de amplitud y los diodos empleados, manteniendo la corriente dentro del rango operativo sin comprometer los componentes.

\end{enumerate}
 
% ══════════════════════════════════════════════════════════════════════════════